\documentclass[../DoAn.tex]{subfiles}
\begin{document}

\section{Mở đầu chương}
Chương này trình bày chi tiết quá trình triển khai Phase 1 của hệ thống RAG v2: thu thập dữ liệu, làm sạch dữ liệu, và xây dựng pipeline chunking. Đây là giai đoạn nền tảng, quyết định chất lượng của toàn bộ hệ thống retrieval ở các phase sau.

Các nội dung triển khai bao gồm:
\begin{itemize}
    \item Thu thập và làm sạch dữ liệu từ 3 nguồn (quydinh, ctdt, stsv);
    \item Xử lý PDF với nhiều công cụ (Docling, PyMuPDF4LLM, olmOCR);
    \item Triển khai chi tiết 4 loại chunker chuyên biệt;
    \item Làm giàu metadata tự động cho chunks CTDT.
\end{itemize}

\section{Thu thập và xử lý dữ liệu quy định (quydinh)}

\subsection{Xử lý PDF với olmOCR}

\subsubsection{Vấn đề font encoding}
Nhiều văn bản PDF quy chế có lỗi font, khi extract text sẽ cho ra ký tự sai:
\begin{verbatim}
Văn bản gốc: "Điều 1. Phạm vi và đối tượng áp dụng"
Lỗi font:    "Äiá»u 1. Phạm vi và đá»i tượng áp dụng"
\end{verbatim}

\subsubsection{Giải pháp olmOCR}
olmOCR sử dụng vision-based OCR, xử lý PDF bằng cách chụp ảnh từng trang rồi nhận diện ký tự bằng mô hình vision, không phụ thuộc vào font embedding:
\begin{verbatim}
# Batch processing 16 files quy dinh
python olmocr/batch_convert.py --input data/quydinh/ 
                               --output data/quydinh/olmocr/
# Kết quả: 16/16 files converted successfully
\end{verbatim}

\subsubsection{So sánh các PDF converter}
\begin{table}[h]
\centering
\begin{tabular}{|l|l|l|l|l|}
\hline
\textbf{Tiêu chí} & \textbf{Docling} & \textbf{PyMuPDF4LLM} & \textbf{olmOCR} \\
\hline
Tốc độ & 8s/34 trang & 3s/34 trang & 60s/34 trang \\
\hline
Độ chính xác TV & 99.5\% & 99.9\% & 99.9\% \\
\hline
Font issues & Có thể fail & Có thể fail & Không ảnh hưởng \\
\hline
Table extraction & Tốt & Trung bình & Xuất sắc \\
\hline
Structure detection & Xuất sắc & Tốt & Tốt \\
\hline
GPU required & Không & Không & Có \\
\hline
\end{tabular}
\caption{So sánh các PDF converter}
\end{table}

\textbf{Quyết định thiết kế:} Sử dụng olmOCR cho văn bản quy định (đảm bảo 100\% success rate), Docling/PyMuPDF4LLM cho CTDT (tốc độ nhanh, structure detection tốt).

\subsection{Làm sạch dữ liệu quy định}
Sau khi OCR, văn bản Markdown cần qua bước làm sạch:
\begin{itemize}
    \item Loại bỏ header/footer lặp lại giữa các trang;
    \item Xóa watermark, page numbers, và các artifact từ OCR;
    \item Chuẩn hóa ký tự đặc biệt tiếng Việt;
    \item Chuyển đổi HTML table (output từ olmOCR) sang Markdown table;
    \item Sửa lỗi ngắt dòng bất thường giữa các đoạn văn.
\end{itemize}

\section{Thu thập và xử lý dữ liệu CTDT}

\subsection{Cấu trúc thư mục}
Mỗi khoa/viện có cấu trúc thư mục chuẩn:
\begin{verbatim}
data/ctdt/
  soict/                    # Viện CNTT&TT
    *.pdf                   # PDF gốc (8 files)
    output_docling/         # Docling OCR output
    clean_data/             # Markdown đã làm sạch
    chunks_recursive.../    # Chunks JSON
  cokhi/                    # Trường Cơ khí
  dien-dientu/              # Viện Điện  
  hoa/                      # Viện KT Hóa học
  toan/                     # Viện Toán UD&TH
  vatlieu/                  # Viện KH&KT Vật liệu
\end{verbatim}

\subsection{Pipeline Docling OCR}
Sử dụng Docling cho các PDF CTDT có layout phức tạp (nhiều bảng, cấu trúc heading):
\begin{verbatim}
# Docling OCR output: markdown với headings
# Ví dụ output:
# Chương trình đào tạo Công nghệ Thông tin
## Thông tin chung
| Ngành đào tạo | Công nghệ Thông tin |
| Mã ngành      | ITE10                |
## Khối kiến thức
### Kiến thức đại cương
| STT | Mã HP | Tên học phần | TC |
|-----|-------|-------------|-----|
| 1   | MI1111| Giải tích I | 4   |
\end{verbatim}

\subsection{Làm sạch dữ liệu CTDT}
Bước làm sạch cho CTDT focus vào:
\begin{itemize}
    \item Chuẩn hóa heading level (H1 cho tiêu đề chính, H2 cho phần, H3 cho mục con);
    \item Sửa bảng markdown bị lỗi format (thiếu separator line, cell không align);
    \item Loại bỏ nội dung trùng lặp giữa các trang;
    \item Thêm suffix \texttt{\_fix.md} cho files đã làm sạch.
\end{itemize}

\section{Thu thập và xử lý dữ liệu STSV}

\subsection{Crawl dữ liệu từ cổng thông tin}
Dữ liệu STSV được crawl từ API của cổng thông tin sinh viên, mỗi tài liệu là một bài viết chứa thông tin hướng dẫn:
\begin{verbatim}
# Ví dụ 1 file JSON:
{
    "DocumentID": 42,
    "Title": "Hướng dẫn sử dụng cổng thông tin sinh viên",
    "TypeDoc": "Sổ tay SV",
    "Description": "## Giới thiệu\n1. Truy cập...\n2. Đăng nhập...",
    "TimeCreate": "2025-08-15T00:00:00"
}
\end{verbatim}

\subsection{Làm sạch dữ liệu STSV}
\textbf{File:} \texttt{data/stsv/clean\_data/clean\_data.py}

Bước làm sạch bao gồm:
\begin{itemize}
    \item Loại bỏ HTML tags còn sót trong Description;
    \item Chuẩn hóa markdown formatting (list, heading, link);
    \item Loại bỏ tài liệu trùng lặp hoặc rỗng;
    \item Gộp tất cả thành file \texttt{data.json} duy nhất để STSVChunker xử lý.
\end{itemize}

\section{Triển khai OlmOcrLegalChunker}

\subsection{Nhận diện cấu trúc văn bản}
Chunker sử dụng các regex patterns để nhận diện cấu trúc pháp lý:
\begin{verbatim}
# Regex patterns for legal structure
CHAPTER_PATTERN = r'^(?:CHƯƠNG|Chương)\s+...'
ARTICLE_PATTERN = r'^(?:Điều|ĐIỀU)\s+(\d+)...'
CLAUSE_PATTERN  = r'^\s*(\d+)\.\s+'
POINT_PATTERN   = r'^\s*[a-zđ]\)\s+'
\end{verbatim}

\subsection{Thuật toán chunking}
\begin{verbatim}
def chunk_document(self, content, source_file):
    # 1. Trích xuất metadata văn bản (căn cứ, ngày ban hành)
    doc_metadata = self._extract_doc_metadata(content)
    
    # 2. Tạo header chunk
    header = self._extract_header(content)
    
    # 3. Chia theo Chương
    chapters = self._split_by_chapters(content)
    
    # 4. Trong mỗi Chương: chia theo Điều
    for chapter in chapters:
        articles = self._split_by_articles(chapter)
        
        for article in articles:
            # 5. Tạo parent chunk (toàn bộ Điều)
            parent = self._create_parent_chunk(article, chapter)
            
            # 6. Tạo child chunks (từng Khoản)
            children = self._create_child_chunks(article, parent)
    
    # 7. Xử lý phụ lục
    appendix = self._extract_appendix(content)
    
    return chunks, stats
\end{verbatim}

\subsection{Readable ID Format}
Mỗi chunk được gán readable ID phản ánh vị trí trong cấu trúc pháp lý:
\begin{verbatim}
def _generate_readable_id(self, chapter, article, clause):
    # c1_d5_k2 = Chương 1, Điều 5, Khoản 2
    parts = []
    if chapter: parts.append(f"c{chapter}")
    if article: parts.append(f"d{article}")
    if clause:  parts.append(f"k{clause}")
    return "_".join(parts)
    
# Ví dụ: c1_d5_k2 → Chương I, Điều 5, Khoản 2
# Ví dụ: c3_d18   → Chương III, Điều 18 (parent, toàn bộ điều)
# Ví dụ: header   → Phần header văn bản
\end{verbatim}

\section{Triển khai RecursiveChunker}

\subsection{Chiến lược H2 Section-based}
RecursiveChunker tách document theo H2 section boundaries:

\begin{verbatim}
def chunk_document(self, text, source=None):
    # 1. Trích xuất tiêu đề document
    doc_title = self._extract_document_title(text)
    
    # 2. Trích xuất tất cả headings
    headings = self._extract_section_headings(text)
    
    # 3. Trích xuất các section H2
    h2_sections = self._extract_h2_sections(text)
    
    # 4. Xử lý phần preamble (trước H2 đầu tiên)
    preamble = text[:h2_sections[0]['start_pos']]
    
    all_chunks = []
    for section in h2_sections:
        # 5. Parent chunk = toàn bộ section H2
        parent = self._create_parent_from_section(section)
        
        # 6. Child chunks = tách section bằng TextSplitter
        children = self._split_text_to_chunks(
            section['content'], doc_title, headings, source
        )
        
        all_chunks.extend([parent] + children)
    
    return all_chunks, stats
\end{verbatim}

\subsection{Bảo vệ bảng Markdown}
Kỹ thuật protect-and-restore để tránh tách giữa bảng:

\begin{verbatim}
def _protect_tables_in_text(self, text):
    """Thay bảng bằng placeholder trước khi split."""
    table_map = {}
    
    def replace_table(match):
        table_text = match.group(0)
        if len(table_text) <= self.chunk_size:
            # Bảng nhỏ: bảo vệ nguyên vẹn
            placeholder = f"__TABLE_{len(table_map):04d}__"
            table_map[placeholder] = table_text
            return placeholder
        # Bảng lớn: để TextSplitter xử lý
        return table_text
    
    protected = re.sub(TABLE_PATTERN, replace_table, text)
    return protected, table_map
\end{verbatim}

Khi chunk chứa bảng lớn cần tách, bảng được chia theo rows với header preserved:
\begin{verbatim}
def _split_table_by_rows(self, table_text, max_rows=0):
    """Tách bảng thành nhiều phần, mỗi phần có header."""
    # header_line + separator_line + N data_rows
    # Tính max_rows dựa trên chunk_size nếu max_rows=0
\end{verbatim}

\subsection{Inject context cho sub-items}
Xử lý trường hợp chunk bắt đầu bằng sub-item không có context:
\begin{verbatim}
# Trước inject:
"b) Đối với sinh viên thuộc diện chính sách..."

# Sau inject (tìm khoản cha từ section text gốc):
"3. Căn cứ kế hoạch tuyển sinh hàng năm:

b) Đối với sinh viên thuộc diện chính sách..."
\end{verbatim}

\section{Triển khai STSVChunker}

\subsection{Phân loại nội dung tự động}
STSVChunker phân tích nội dung Description và chọn chiến lược phù hợp:

\begin{verbatim}
def _segment(self, description, title, type_doc):
    # Loại 1: Ngắn (< 1500 chars) -> 1 chunk
    if len(description) <= SINGLE_CHUNK_THRESHOLD:
        return [(description, None, None)]
    
    # Loại 2: Có phần Roman (I., II., III.)
    if _RE_ROMAN.search(description):
        return self._split_by_roman_then_items(description)
    
    # Loại 3: Có mục số (1., 2., ...)
    if _RE_NUMBERED.search(description):
        return self._split_by_numbered_items(description)
    
    # Loại 4: Văn xuôi thuần
    return [(description, None, None)]
\end{verbatim}

\subsection{Sequential tracking cho mục số}
Để tránh nhầm lẫn sub-items (ví dụ ``1. Miễn học phần...'' xuất hiện trong mục 8) với mục cấp cao, STSVChunker sử dụng sequential tracking:
\begin{verbatim}
# Chỉ nhận mục mới khi số thứ tự đúng = expected
# 1 -> 2 -> 3 -> ... (tuần tự)
expected_num = None
for line in lines:
    m = re.match(r'^(\d{1,2})\.\s+\S', stripped)
    if m:
        num = int(m.group(1))
        if expected_num is None and num == 1:
            expected_num = 1
        if expected_num is not None and num == expected_num:
            # Đây là mục cấp cao mới
            expected_num = num + 1
\end{verbatim}

\subsection{Context prefix cho retrieval}
Mỗi chunk được thêm prefix để đảm bảo context khi retrieval:
\begin{verbatim}
def _build_content(self, text, title, type_doc, section_ctx):
    # Prefix: [Title | TypeDoc | Section_context]
    parts = [title]
    if type_doc:     parts.append(type_doc)
    if section_ctx:  parts.append(section_ctx)
    prefix = " | ".join(parts)
    return f"[{prefix}]\n{text}"

# Ví dụ output:
# [Hướng dẫn đăng ký học tập | Sổ tay SV | II. Đăng ký học phần]
# 1. Truy cập hệ thống SIS tại https://sis.hust.edu.vn
# 2. Chọn mục "Đăng ký học phần"
\end{verbatim}

\section{Triển khai Metadata Enrichment cho CTDT}

\subsection{Trích xuất ngày ban hành}
Thuật toán trích xuất ngày ban hành theo thứ tự ưu tiên:
\begin{verbatim}
def extract_effective_date(text, filename):
    # Priority 1: Date in filename (YYYY.MM.DD, YYYY-MM-DD)
    # Priority 2: "ngày DD tháng MM năm YYYY" trong 2000 char đầu
    # Priority 3: "ngày DD/MM/YYYY" trong 2000 char đầu
    # Priority 4: Năm trong tiêu đề "CHƯƠNG TRÌNH ... 2017"
    # Priority 5: Năm trong filename
\end{verbatim}

\subsection{Xác định khóa áp dụng}
Phát hiện khóa áp dụng từ filename và nội dung:
\begin{verbatim}
def extract_applicable_cohort(text, filename):
    # Tìm pattern K65, K68, K70, ... trong filename và content
    # Trả về danh sách sorted: "K65, K68, K70"
\end{verbatim}

\subsection{Xác định ngành đào tạo}
Chiến lược nhiều tầng để xác định ngành:
\begin{enumerate}
    \item Trường ``Ngành đào tạo:'' trong nội dung (plain text hoặc table);
    \item Trường ``Tên chương trình:'' trong nội dung;
    \item H1 heading nếu là tên ngành (không phải heading tổng quát);
    \item Program code mapping từ filename (IT2 $\rightarrow$ ``Kỹ thuật Máy tính'');
    \item Keyword patterns trong filename;
    \item H2 heading đầu tiên nếu là tên ngành.
\end{enumerate}

\subsection{Phân loại loại tài liệu}
Phân loại tự động dựa trên filename và tiêu đề:
\begin{verbatim}
DOCUMENT_TYPE_RULES = [
    (r"tài\s*năng|talent",       "talent_program"),
    (r"tiên\s*tiến|CTTT",        "advanced_program"),
    (r"chất\s*lượng\s*cao|KSCLC","high_quality_program"),
    (r"việt.*pháp|ITEP",         "international_program"),
    (r"tích\s*hợp.*cử\s*nhân",  "integrated_program"),
    (r"khung\s*CT",              "training_framework"),
]
# Default: "curriculum"
\end{verbatim}

%% =====================================================================
%% PHASE 4: Tool Search (Tavily) + MongoDB Memory
%% =====================================================================

\section{Triển khai Phase 4: Tool Search và MongoDB Memory}

Phase 4 bổ sung hai thành phần quan trọng: (1) Tavily Web Search -- cơ chế fallback khi RAG không đủ context để trả lời, và (2) MongoDB Memory -- tầng lưu trữ lịch sử hội thoại và trạng thái phiên làm việc.

\subsection{Tavily Search Tool}
\textbf{File:} \texttt{tools/tavily\_search.py} (145 dòng, 5.793 bytes)

\subsubsection{Vai trò trong hệ thống}
Tavily Search Tool đóng vai trò fallback khi Self Evaluator đánh giá câu trả lời từ RAG không đạt yêu cầu. Khi đó, hệ thống tự động tìm kiếm web để bổ sung context mới, sau đó sinh lại câu trả lời:
\begin{verbatim}
SelfEvaluator → FAIL
  → TavilySearchTool.search(query)
  → Lấy web context mới
  → ChatModel.generate(query, web_context)
  → Final Answer (sinh lại từ web context)
\end{verbatim}

\subsubsection{Kiến trúc class TavilySearchTool}
\begin{verbatim}
class TavilySearchTool:
    def __init__(self, api_key, max_results=5,
                 max_retries=3, min_retry_delay=1.0):
        self._client = TavilyClient(api_key=resolved_key)
        self._last_call_time = 0.0  # Rate limiting

    def search(self, query, max_results=None,
               search_depth="basic", include_answer=True):
        # Rate limiting + Exponential backoff retry
        # Returns: {query, answer, results, context}

    def _parse_results(response) -> List[Dict]:
        # Extract: title, url, content từ raw response

    def _format_context(results) -> str:
        # Format thành numbered text block cho LLM
\end{verbatim}

\subsubsection{Các tính năng chính}
\begin{itemize}
    \item \textbf{Rate limiting}: Đảm bảo khoảng cách tối thiểu 1 giây giữa các API call, tránh bị giới hạn bởi Tavily;
    \item \textbf{Exponential backoff}: Retry tối đa 3 lần với delay tăng dần ($1s \rightarrow 2s \rightarrow 4s$) khi gặp lỗi tạm thời;
    \item \textbf{Format context cho LLM}: Kết quả web search được format thành chuỗi có đánh số, phù hợp làm context cho prompt RAG:
\end{itemize}

\begin{verbatim}
# Output format của _format_context():
[1] Tiêu đề bài viết 1
URL: https://example.com/article1
Nội dung trích xuất từ trang web...

---

[2] Tiêu đề bài viết 2
URL: https://example.com/article2
Nội dung trích xuất từ trang web...
\end{verbatim}

\subsubsection{Cấu hình API}
\begin{verbatim}
TavilySearchTool(
    api_key="TAVILY_API_KEY",     # Từ .env
    max_results=5,                # Số kết quả mặc định
    max_retries=3,                # Số lần retry
    min_retry_delay=1.0,          # Delay cơ sở (giây)
)
# search_depth: "basic" (nhanh) hoặc "advanced" (sâu hơn)
# include_answer: Tavily tự sinh short answer
\end{verbatim}

\subsection{MongoDB Memory Layer}

Tầng MongoDB Memory gồm 3 module: \texttt{MongoClient} (kết nối), \texttt{ChatHistoryStore} (lịch sử chat), và \texttt{ConversationState} (trạng thái phiên).

\subsubsection{MongoClient -- Quản lý kết nối}
\textbf{File:} \texttt{memory/mongo\_client.py} (66 dòng, 2.376 bytes)

MongoClient là wrapper mỏng quanh \texttt{pymongo.MongoClient} với connection pooling:
\begin{verbatim}
class MongoClient:
    def __init__(self, uri="mongodb://localhost:27017",
                 database="rag_chatbot", max_pool_size=10):
        self._client = PyMongoClient(
            uri,
            maxPoolSize=max_pool_size,
            serverSelectionTimeoutMS=5000,
        )
        self._db = self._client[database]

    @property
    def db(self) -> Database:
        """Return the active database handle."""

    def ping(self) -> bool:
        """Check connectivity."""

    def close(self) -> None:
        """Close the connection pool."""
\end{verbatim}

\textbf{Quyết định thiết kế:}
\begin{itemize}
    \item \textbf{Connection pooling} (\texttt{maxPoolSize=10}): Tái sử dụng kết nối, giảm overhead cho các request liên tục;
    \item \textbf{Server selection timeout} (5 giây): Fail fast khi MongoDB không khả dụng, tránh block pipeline;
    \item Database mặc định: \texttt{rag\_chatbot}.
\end{itemize}

\subsubsection{ChatHistoryStore -- Lưu trữ lịch sử chat}
\textbf{File:} \texttt{memory/chat\_history.py} (128 dòng, 4.226 bytes)

ChatHistoryStore cung cấp CRUD operations cho tin nhắn chat, lưu vào collection \texttt{chat\_messages}:
\begin{verbatim}
class ChatHistoryStore:
    def __init__(self, mongo_client, collection_name="chat_messages"):
        self._collection = mongo_client.db[collection_name]
        self._ensure_indexes()  # Index: (session_id, timestamp)

    def save_message(self, session_id, role, content) -> str:
        """Lưu 1 tin nhắn. Returns: inserted _id."""
        doc = {
            "session_id": session_id,
            "role": role,          # "user" hoặc "assistant"
            "content": content,
            "timestamp": datetime.now(timezone.utc),
        }

    def get_history(self, session_id, limit=20) -> List[Dict]:
        """Lấy N tin nhắn gần nhất, sorted oldest -> newest."""
        # Sort by timestamp DESC, limit, rồi reverse

    def clear_history(self, session_id) -> int:
        """Xóa toàn bộ lịch sử của 1 session."""
\end{verbatim}

\textbf{Index strategy:} Compound index \texttt{(session\_id, timestamp)} cho phép truy vấn nhanh theo session và sắp xếp theo thời gian mà không cần full collection scan.

\subsubsection{ConversationState -- Quản lý trạng thái phiên}
\textbf{File:} \texttt{memory/conversation.py} (163 dòng, 5.684 bytes)

ConversationState quản lý vòng đời phiên hội thoại và lưu trữ câu trả lời cuối cùng kèm metadata:
\begin{verbatim}
class ConversationState:
    def __init__(self, mongo_client):
        self._conversations = mongo_client.db["conversations"]
        self._answers = mongo_client.db["answers"]
        self._ensure_indexes()

    # Session lifecycle
    def create_session(self, session_id):
        """Upsert session: tạo mới hoặc touch nếu đã tồn tại."""

    def update_status(self, session_id, status):
        """Cập nhật status: "active" -> "closed"."""

    def touch(self, session_id):
        """Bump last_active timestamp."""

    def get_session(self, session_id) -> Optional[Dict]:
        """Trả về metadata phiên."""

    # Final answers
    def save_answer(self, session_id, query, answer,
                    sources=None, scores=None) -> str:
        """Lưu câu trả lời + provenance (sources, scores)."""

    def get_answers(self, session_id, limit=10) -> List[Dict]:
        """Lấy các câu trả lời gần nhất."""
\end{verbatim}

\textbf{Cấu trúc dữ liệu MongoDB:}
\begin{verbatim}
# Collection: conversations
{
    "session_id": "uuid-string",
    "status": "active" | "closed",
    "created_at": ISODate("2026-03-29T..."),
    "last_active": ISODate("2026-03-29T...")
}

# Collection: answers
{
    "session_id": "uuid-string",
    "query": "Điều kiện tốt nghiệp là gì?",
    "answer": "Theo Điều 5, sinh viên cần...",
    "sources": [{"title": "...", "score": 0.95}],
    "scores": {"relevance": 0.9, "faithfulness": 0.85},
    "timestamp": ISODate("2026-03-29T...")
}
\end{verbatim}

\textbf{Quyết định thiết kế:}
\begin{itemize}
    \item \textbf{Tách conversations và answers}: Cho phép truy vấn analytics độc lập (ví dụ: thống kê chất lượng câu trả lời theo thời gian);
    \item \textbf{Upsert pattern} cho \texttt{create\_session}: Idempotent -- gọi nhiều lần không tạo duplicate;
    \item \textbf{Lưu sources + scores}: Traceability đầy đủ, biết câu trả lời dựa trên tài liệu nào và chất lượng ra sao.
\end{itemize}

\subsection{MongoLogger -- Logging tích hợp cho Pipeline}
\textbf{File:} \texttt{pipeline/mongo\_logger.py} (149 dòng, 5.316 bytes)

MongoLogger là module logging cấp cao hơn, được tích hợp trực tiếp vào RAG Pipeline để ghi log mỗi lượt hỏi-đáp:
\begin{verbatim}
class MongoLogger:
    def __init__(self, uri, database):
        self._sessions = self._db["sessions"]
        self._query_logs = self._db["query_logs"]

    def new_session(self, user_id=None) -> str:
        """Tạo session mới với UUID, returns session_id."""

    def log_turn(self, session_id, question, result,
                 reflected_question=None, latency_ms=0) -> int:
        """Ghi 1 lượt hỏi-đáp vào session + query_logs."""

    def get_history(self, session_id, max_turns=10):
        """Trả về lịch sử dạng [{role, content}]."""

    def get_session(self, session_id) -> Optional[Dict]:
        """Trả về full session document."""
\end{verbatim}

\textbf{Cấu trúc session document:}
\begin{verbatim}
# Collection: sessions
{
    "session_id": "uuid",
    "user_id": "optional",
    "created_at": ISODate(...),
    "updated_at": ISODate(...),
    "turn_count": 3,
    "turns": [
        {
            "turn_id": 1,
            "question": "Điều kiện tốt nghiệp?",
            "answer": "Theo quy chế...",
            "intent": "rag",
            "num_sources": 5,
            "latency_ms": 1200,
            "timestamp": ISODate(...)
        },
        ...
    ]
}

# Collection: query_logs (flat, cho analytics)
{
    "session_id": "uuid",
    "turn_id": 1,
    "question": "...",
    "answer": "...",
    "intent": "rag",
    "model_name": "gemini-2.5-flash",
    "latency_ms": 1200,
    "timestamp": ISODate(...)
}
\end{verbatim}

\textbf{Quyết định thiết kế:}
\begin{itemize}
    \item \textbf{Dual storage}: Embedded turns trong session (đọc nhanh toàn bộ hội thoại) + flat query\_logs (analytics, thống kê);
    \item \textbf{Tự động tính turn\_id}: Dựa trên \texttt{turn\_count} hiện tại, đảm bảo thứ tự tuần tự;
    \item \textbf{Ghi latency}: Đo thời gian xử lý mỗi lượt, phục vụ monitoring và tối ưu.
\end{itemize}

%% =====================================================================
%% PHASE 5: FastAPI Backend + Pipeline Integration
%% =====================================================================

\section{Triển khai Phase 5: FastAPI Backend và Pipeline Integration}

Phase 5 kết nối tất cả các tầng đã triển khai (Phase 1--4) thành một pipeline thống nhất, đồng thời cung cấp API server qua FastAPI với streaming response.

\subsection{Pipeline Orchestration}

\subsubsection{RAGPipeline -- Điều phối tổng thể}
\textbf{File:} \texttt{pipeline/rag\_pipeline.py} (327 dòng, 12.283 bytes)

RAGPipeline là class trung tâm, khởi tạo và kết nối tất cả components:
\begin{verbatim}
class RAGPipeline:
    def __init__(self, settings=None, api_key=None,
                 config=None, mongo_logger=None):
        # 1. Load settings từ .env
        settings = Settings()

        # 2. Khởi tạo Query Router (classifier-based)
        self._router = QueryRouter(mode="classifier")

        # 3. Khởi tạo Query Reflector (LLM-based rewrite)
        self._reflector = QueryReflector()

        # 4. Load 2 embedding models
        self._bge = BGEm3Embedder()        # BAAI/bge-m3
        self._e5 = E5MultilingualEmbedder() # multilingual-e5-large

        # 5. Kết nối retrieval stores (Qdrant + Elasticsearch)
        self._searcher = create_retriever(settings)

        # 6. Load reranker
        self._reranker = create_reranker(settings)

        # 7. Khởi tạo Chat model
        self._chat = create_llm(settings)   # Gemini 2.5 Flash

        # 8. Self evaluator (tái sử dụng LLM instance)
        self._self_eval = SelfEvaluator(llm=self._chat)

        # 9. Tavily fallback
        self._tavily = TavilySearchTool(api_key=tavily_key)
\end{verbatim}

\textbf{Quyết định thiết kế quan trọng:}
\begin{itemize}
    \item \textbf{Singleton pattern}: Tất cả models được load 1 lần duy nhất khi khởi tạo pipeline, tránh reload giữa các request;
    \item \textbf{Factory pattern}: Sử dụng \texttt{create\_llm()}, \texttt{create\_retriever()}, \texttt{create\_reranker()} -- cho phép đổi provider qua \texttt{.env} mà không sửa code;
    \item \textbf{Self-eval tái sử dụng LLM}: \texttt{SelfEvaluator} dùng cùng instance \texttt{BaseLLM} với chat model, không tạo thêm API client.
\end{itemize}

\subsubsection{Public API}
RAGPipeline cung cấp 2 phương thức chính:
\begin{verbatim}
def query(self, question, history=None, top_k=None,
          session_id=None) -> Dict:
    """Non-streaming: trả về full answer một lần."""
    # 1. Auto-load history từ MongoDB nếu có session_id
    # 2. Route query (chitchat/rag)
    # 3. Gọi flow tương ứng
    # 4. Log kết quả vào MongoDB
    # Returns: {question, answer, sources, num_sources,
    #           intent, model_name}

def query_stream(self, question, history=None, top_k=None,
                 session_id=None) -> Generator[str]:
    """Streaming: yield từng token."""
    # Tương tự query(), nhưng generation streaming
    # Log vào MongoDB sau khi stream kết thúc
\end{verbatim}

\subsubsection{Pipeline Flows}
\textbf{File:} \texttt{pipeline/flows.py} (304 dòng, 10.896 bytes)

Module flows định nghĩa 2 luồng xử lý tách biệt:

\textbf{Chitchat Flow:}
\begin{verbatim}
def chitchat_flow(question, history, chat_model) -> Dict:
    """Router -> Chat Model -> Response (không retrieval)."""
    trimmed = _trim_history(history, limit=6)
    answer = chat_model.generate(
        query=question, history=trimmed, mode="chitchat"
    )
    return {"answer": answer, "sources": [], "intent": "chitchat"}
\end{verbatim}

\textbf{RAG Flow (luồng chính):}
\begin{verbatim}
def rag_flow(question, history, reflector, bge_embedder,
             e5_embedder, searcher, reranker, chat_model,
             self_evaluator, tavily_tool, cfg) -> Dict:
    # 1. Reflection: rewrite query cho retrieval tốt hơn
    search_query = reflector.reflect(question, chat_history)

    # 2. Embed: tạo vector cho cả BGE-M3 và E5
    bge_vec = bge_embedder.embed_query(search_query)
    e5_vec = e5_embedder.embed_query(search_query)

    # 3. Hybrid Search: Qdrant + Elasticsearch + RRF fusion
    raw_results = searcher.search(
        query=search_query,
        bge_m3_query=bge_vec, e5_query=e5_vec,
        top_k=cfg["top_k"] * 4,  # Pool 4x để reranker chọn
    )

    # 4. Rerank: chọn top-K tài liệu chính xác nhất
    reranked = reranker.rerank(
        query=search_query, documents=raw_results,
        top_k=cfg["top_k"]   # Mặc định: 5
    )

    # 5. Generate: LLM sinh câu trả lời từ context
    context = _format_context(reranked)
    answer = chat_model.generate(
        query=question, context=context, mode="rag"
    )

    # 6. Self-eval: kiểm tra chất lượng
    eval_result = self_evaluator.evaluate(
        query=question, context=context, response=answer
    )
    if not eval_result["pass"]:
        # 7. Tavily fallback: tìm web, sinh lại
        answer = _tavily_fallback(question, tavily_tool,
                                  chat_model, history)

    return {"answer": answer, "sources": reranked, ...}
\end{verbatim}

\textbf{Tavily Fallback:}
\begin{verbatim}
def _tavily_fallback(question, answer, tavily_tool,
                     chat_model, history) -> str:
    """Web search -> re-generate khi self-eval FAIL."""
    search_result = tavily_tool.search(question)
    web_context = search_result["context"]
    new_answer = chat_model.generate(
        query=question, context=web_context, mode="rag"
    )
    return new_answer
\end{verbatim}

\textbf{Helper functions:}
\begin{itemize}
    \item \texttt{\_format\_context(documents)}: Chuyển danh sách tài liệu retrieved thành chuỗi context có đánh số, mỗi tài liệu gồm title và text;
    \item \texttt{\_trim\_history(history, limit=6)}: Giới hạn lịch sử chat gửi cho LLM, giữ 6 lượt gần nhất để tiết kiệm token.
\end{itemize}

\subsection{FastAPI Backend}

\subsubsection{Application Entry Point}
\textbf{File:} \texttt{api/main.py} (117 dòng, 3.491 bytes)

Ứng dụng FastAPI sử dụng lifespan pattern để khởi tạo pipeline 1 lần duy nhất:
\begin{verbatim}
@asynccontextmanager
async def lifespan(app: FastAPI):
    """Startup: load pipeline (models, connections)."""
    settings = Settings()

    # 1. Khởi tạo MongoLogger (nếu MongoDB enabled)
    mongo_logger = MongoLogger(
        uri=settings.mongodb_uri,
        database=settings.mongodb_database
    )
    app.state.mongo_logger = mongo_logger

    # 2. Load RAG Pipeline (blocking, chạy trong executor)
    app.state.pipeline = await loop.run_in_executor(
        None,
        lambda: RAGPipeline(api_key=api_key,
                            mongo_logger=mongo_logger),
    )

    yield  # App is running
    # Shutdown cleanup

def create_app() -> FastAPI:
    app = FastAPI(
        title="RAG v2 Chatbot API",
        version="2.0.0",
        lifespan=lifespan,
    )

    # CORS middleware
    app.add_middleware(CORSMiddleware,
        allow_origins=["*"],
        allow_methods=["GET","POST","PUT","DELETE","OPTIONS"],
    )

    # Register routes
    app.include_router(chat_router)
    app.include_router(health_router)
    app.include_router(session_router)

    return app
\end{verbatim}

\textbf{Quyết định thiết kế:}
\begin{itemize}
    \item \textbf{Lifespan pattern} (thay vì \texttt{on\_startup}): Chuẩn mới của FastAPI, đảm bảo cleanup khi shutdown;
    \item \textbf{run\_in\_executor}: Load models trong thread pool, không block event loop;
    \item \textbf{CORS wildcard}: Cho phép frontend từ mọi origin kết nối (phù hợp giai đoạn phát triển);
    \item \textbf{State pattern}: Pipeline và MongoLogger lưu trong \texttt{app.state}, truy cập từ mọi route.
\end{itemize}

\subsubsection{Chat Endpoints}
\textbf{File:} \texttt{api/routes/chat.py} (139 dòng, 4.574 bytes)

Hai endpoint chính cho chat:

\textbf{POST /chat} -- Non-streaming response:
\begin{verbatim}
@router.post("/chat", response_model=ChatResponse)
async def chat(request: Request, body: ChatRequest):
    pipeline = request.app.state.pipeline
    mongo_logger = request.app.state.mongo_logger

    # Auto-create session nếu chưa có
    if body.session_id is None and mongo_logger:
        session_id = mongo_logger.new_session()

    # Chạy pipeline trong thread pool (non-blocking)
    result = await loop.run_in_executor(
        None,
        lambda: pipeline.query(
            question=body.question,
            history=history,
            top_k=body.top_k,
            session_id=session_id,
        ),
    )

    return ChatResponse(
        question=result["question"],
        answer=result["answer"],
        retrieved_documents=retrieved_docs,
        num_documents=result["num_sources"],
        model_name=result["model_name"],
        intent=result["intent"],
        session_id=session_id,
    )
\end{verbatim}

\textbf{POST /chat/stream} -- Server-Sent Events (SSE):
\begin{verbatim}
@router.post("/chat/stream")
async def chat_stream(request, body) -> StreamingResponse:
    async def event_generator():
        # Gửi session_id như SSE event đầu tiên
        yield f'data: {{"session_id": "{session_id}"}}\n\n'

        # Producer-consumer pattern với asyncio.Queue
        queue = asyncio.Queue()

        def _produce():
            for chunk in pipeline.query_stream(...):
                loop.call_soon_threadsafe(
                    queue.put_nowait, chunk)
            loop.call_soon_threadsafe(
                queue.put_nowait, None)  # sentinel

        loop.run_in_executor(None, _produce)

        while True:
            chunk = await queue.get()
            if chunk is None:
                yield "data: [DONE]\n\n"
                break
            yield f"data: {chunk}\n\n"

    return StreamingResponse(
        event_generator(),
        media_type="text/event-stream"
    )
\end{verbatim}

\textbf{Quyết định thiết kế SSE streaming:}
\begin{itemize}
    \item \textbf{Producer-consumer với asyncio.Queue}: Pipeline chạy synchronous trong thread pool (producer), SSE event loop đọc từ queue (consumer) -- bridge giữa sync và async;
    \item \textbf{Session ID gửi trước}: Frontend nhận session\_id ngay lập tức, không cần đợi câu trả lời;
    \item \textbf{Sentinel \texttt{[DONE]}}: Chuẩn SSE signaling, frontend biết stream đã kết thúc.
\end{itemize}

\subsubsection{Health Check và Session Endpoints}
\textbf{File:} \texttt{api/routes/health.py} (19 dòng), \texttt{api/routes/session.py} (51 dòng)

\begin{verbatim}
# GET /health
@router.get("/health", response_model=HealthResponse)
async def health(request):
    return HealthResponse(status="healthy", rag_initialized=True)

# POST /session — tạo session mới
@router.post("")
async def create_session(request):
    session_id = mongo_logger.new_session()
    return {"session_id": session_id, "created_at": ...}

# GET /session/{session_id} — lấy chi tiết session
@router.get("/{session_id}")
async def get_session(request, session_id):
    return mongo_logger.get_session(session_id)
\end{verbatim}

\subsubsection{API Schemas (Pydantic)}
\textbf{File:} \texttt{api/schemas.py} (62 dòng, 2.213 bytes)

Sử dụng Pydantic BaseModel cho validation và serialization:
\begin{verbatim}
# Request
class ChatRequest(BaseModel):
    question: str = Field(..., min_length=1, max_length=4096)
    top_k: int = Field(default=5, ge=1, le=50)
    history: Optional[List[HistoryMessage]] = None
    session_id: Optional[str] = None

class HistoryMessage(BaseModel):
    role: str = Field(..., pattern="^(user|assistant)$")
    content: str

# Response
class ChatResponse(BaseModel):
    question: str
    answer: str
    retrieved_documents: List[RetrievedDocument]
    num_documents: int
    model_name: str
    intent: str
    session_id: str

class RetrievedDocument(BaseModel):
    rank: int
    content: str
    score: float
    metadata: Dict[str, Any]
\end{verbatim}

\textbf{Validation rules:}
\begin{itemize}
    \item \texttt{question}: Bắt buộc, 1--4096 ký tự;
    \item \texttt{top\_k}: 1--50, mặc định 5;
    \item \texttt{role}: Chỉ chấp nhận \texttt{"user"} hoặc \texttt{"assistant"} (regex pattern).
\end{itemize}

\subsection{Centralized Configuration}

\subsubsection{Settings class}
\textbf{File:} \texttt{config/settings.py} (133 dòng, 4.898 bytes)

Toàn bộ cấu hình hệ thống được quản lý tập trung qua Pydantic BaseSettings:
\begin{verbatim}
class Settings(BaseSettings):
    # Provider selectors (đổi qua .env, không sửa code)
    llm_provider: str = "gemini"     # gemini|openai|azure|ollama
    embedding_provider: str = "ensemble" # ensemble|bge_m3|e5
    reranker_provider: str = "bge"    # bge|cohere|none

    # API Keys
    google_api_key: str = ""
    openai_api_key: str = ""
    tavily_api_key: str = ""

    # Infrastructure
    qdrant_host: str = "localhost"
    qdrant_port: int = 6333
    elasticsearch_host: str = "localhost"
    elasticsearch_port: int = 9200
    mongodb_uri: str = "mongodb://localhost:27017"
    mongodb_database: str = "rag_chatbot"
    mongodb_enabled: bool = True

    # Collections
    collections: List[str] = ["stsv","quydinh","kehoach","ctdt"]

    # Chat Model
    chat_model: str = "gemini-2.5-flash"
    chat_temperature: float = 0.3
    chat_max_tokens: int = 5120

    # Retrieval
    top_k: int = 5
    vector_weight: float = 0.8
    keyword_weight: float = 0.2

    # Reranker
    reranker_model: str = "BAAI/bge-reranker-v2-m3"
    reranker_top_k: int = 5

    # Features
    self_eval_enabled: bool = True
    tavily_fallback_enabled: bool = True
    reflection_enabled: bool = True

    # Server
    api_host: str = "0.0.0.0"
    api_port: int = 8000

    model_config = {
        "env_file": ".env",
        "extra": "ignore",
    }
\end{verbatim}

\textbf{Ưu điểm của Pydantic BaseSettings:}
\begin{itemize}
    \item \textbf{Tự động load từ .env}: Không cần parse thủ công;
    \item \textbf{Type validation}: Đảm bảo kiểu dữ liệu đúng (int, float, List);
    \item \textbf{Default values}: Mọi config đều có giá trị mặc định hợp lý;
    \item \textbf{Override qua env var}: Deploy trên server chỉ cần set env var, không sửa code;
    \item \textbf{Provider pattern}: Đổi LLM từ Gemini sang OpenAI chỉ cần set \texttt{LLM\_PROVIDER=openai} trong \texttt{.env}.
\end{itemize}

\subsection{Tổng kết API Endpoints}

\begin{table}[H]
\centering
\begin{tabular}{|l|l|p{7cm}|}
\hline
\textbf{Method} & \textbf{Endpoint} & \textbf{Mô tả} \\
\hline
GET & /health & Health check, kiểm tra pipeline đã sẵn sàng \\
\hline
POST & /chat & Chat non-streaming, trả đầy đủ answer + sources \\
\hline
POST & /chat/stream & Chat streaming SSE, trả từng token \\
\hline
POST & /session & Tạo session mới \\
\hline
GET & /session/\{id\} & Lấy chi tiết session (gồm tất cả turns) \\
\hline
\end{tabular}
\caption{Danh sách API endpoints}
\end{table}

\section{Kết luận chương}
Chương này đã trình bày chi tiết quá trình triển khai hệ thống RAG v2 qua các phase:

\begin{itemize}
    \item \textbf{Phase 1 -- Thu thập và xử lý dữ liệu}: 16 PDF quy định, CTDT từ 6 khoa/viện, 83 tài liệu JSON STSV; 4 loại chunker chuyên biệt tạo ra 2.752 chunks chất lượng cao;
    \item \textbf{Phase 4 -- Tool Search và MongoDB Memory}: Tavily web search với rate limiting và exponential backoff làm fallback; MongoDB lưu trữ lịch sử chat, trạng thái phiên, và câu trả lời kèm provenance;
    \item \textbf{Phase 5 -- FastAPI Backend và Pipeline Integration}: RAGPipeline kết nối 8 tầng thành pipeline end-to-end; FastAPI server với SSE streaming, health check, và session management; Pydantic BaseSettings cho centralized configuration.
\end{itemize}

Kết quả, hệ thống sẵn sàng nhận request từ frontend, xử lý qua toàn bộ pipeline (Router $\rightarrow$ Reflection $\rightarrow$ Embed $\rightarrow$ Search $\rightarrow$ Rerank $\rightarrow$ Generate $\rightarrow$ Self-Eval $\rightarrow$ Fallback), và trả kết quả streaming cho người dùng.

\end{document}
